% SHARD-ID: CA-74-JW-KERNEL-DECISION
% SHARD-TITLE: The Jones-Wenzl Kernel Decision for the Dilute Evaluation Map
% SHARD-SUMMARY: Settles the CA-69 conjecture at L=4: the dilute evaluation map rho_4 surjects onto the parity-preserving subalgebra and its one-dimensional kernel is the corner-embedded Jones-Wenzl projector iota(p_4).
% SHARD-SUMMARY: Derives the parity-refined charge multiplicities in closed form, (F_{2L-1} +- F_{L+1})/2 and (F_{2L} -+ F_L)/2, giving parity-preserving dimensions 9, 51, 322, 2135, 14445 for L = 2..6.
% SHARD-SUMMARY: Upgrades CA-69 to the conjecture ker rho_L = the ideal generated by corner Jones-Wenzl embeddings, with rigorous lower bounds 53 and 1066 on the kernel dimension at L = 5, 6.
% SHARD-KEYWORDS: Jones-Wenzl, kernel, dilute Temperley-Lieb, parity, evaluation map, negligible ideal, radical, Fibonacci, rank

\section{The Jones--Wenzl Kernel Decision for the Dilute Evaluation Map}
\label{sec:jw-kernel-decision}

\subsection{Status, Sources, and Dependencies}

\textbf{Status: derived, with one numerical input.}  This shard interprets
the CA-73 decision data and settles the CA-69 conjecture at \(L=4\).  The
combinatorics and the kernel identification are local derivations from
registered sources; the single essential numerical input is the computed
rank \(\operatorname{rank}\rho_4=322\) (CA-73, re-run here).  Executes item
W1.4 of \texttt{reviews/2026-07-06\_pipeline\_consistency\_scoping/PLAN.md};
orchestration record in
\texttt{reviews/2026-07-06\_ca74\_jw\_kernel\_decision/}.

Dependencies: CA-66, CA-69, CA-71, CA-72, CA-73; CONVENTIONS (b), (r).

% Source: references/text/TemperleyLiebRootsJonesQuotient.txt:63--69 --
%   (Iohara-Lehrer-Zhang, SRC-TL-JONES) at |q^2| = l the Jones trace form on
%   TL_{l-1} has a radical of dimension 1 generated by the Jones-Wenzl
%   idempotent; for n >= l-1 the radical R_n is generated by E_{l-1}; the
%   Jones algebra Q_n(l) = TL_n/R_n is the canonical semisimple quotient.
% Source: references/text/TemperleyLiebRootsJonesQuotient.txt:74--80 --
%   Q_n(l) = End_{U_q(sl2)}(Delta_q(1)^{tensor n}) in the reduced tilting
%   (fusion) category of Reshetikhin-Turaev-Andersen.
% Source: references/text/TemperleyLiebRootsJonesQuotient.txt:135--141 --
%   Prop 3.1: unique idempotent E_{l-1} in TL_{l-1} with f_i E = E f_i = 0;
%   trace non-degenerate iff n <= l-2; n = l-1 degenerate yet semisimple.
% Source: references/text/TemperleyLiebRootsJonesQuotient.txt:792 --
%   "the coefficient of I^{tensor(l-1)} in E_{l-1} is 1".
% Source: references/text/TemperleyLiebRootsJonesQuotient.txt:818 --
%   "tr_{l-1}(E_{l-1}) = 0" (negligibility of the first JW at the root).
% Source: references/lattice-symmetry/DiluteTL2014/source/Dtl_pgl_ell.06.tex:891--895
%   -- corner iso pi_A dTL_n pi_A = TL_|A|; :2410--2415 -- the same via
%   pi_z, and S_n = direct sum of corners; :773--776 -- parity ideals;
%   :1020--1033 -- dim dTL_n = M_{2n}.
% Source: references/lattice-symmetry/DiluteTL2014/source/Dtl_pgl_ell.06.tex:2314--2318
%   -- the dilute radical decomposes as a multiplicity-binomial direct sum
%   of dense TL Gram radicals (used for the L >= 5 outlook only).
% Source: literature/md/1710.07362/1710.07362.md:30--33 -- Wenzl recursion
%   f^{(k+1)} = f^{(k)} - ([k]_q/[k+1]_q)(...); the extraction truncates the
%   diagram factor of the second term, so the Julia object is pinned by its
%   checked defining properties, not by this formula alone (see below).
% Source: report/sections/66_anyonic_states_variable_n_fock.tex:194--205 --
%   m_1(L) = F_{2L-1}, m_tau(L) = F_{2L}, dim A_L = F_{4L-1}.
% Source: report/sections/69_dilute_tl_word_algebra.tex:70--128 -- the
%   evaluation map rho_L, the parity obstruction, and the conjecture this
%   shard settles at L=4.
% Source: worklog/013_2026-07-06_categorical_residuals.md:46--63 -- the
%   CA-73 decision data 322/322/323 (independently re-run).
% Source: src/JonesWenzlKernel.jl and test/runtests.jl testset
%   "jones-wenzl kernel decision (CA-74)" -- the checked shadows below.

\subsection{Setting and the Decision Data}

Fix Fibonacci \((\Cc,O)\), \(O=\one\oplus\tau\), per CONVENTIONS (r), the
tower \(\Alg_L=\operatorname{End}_{\Cc}(O^{\otimes L})\), and the CA-69
evaluation map \(\rho_L:\mathrm{dTL}_L(\varphi)\to\Alg_L\).  CA-73 computed
(and this session re-ran) the decision data: the image of the Tier-1 dilute
generators has dimension \(9,\,51,\,322\) at \(L=2,3,4\), equal at every
\(L\) to the dimension of the parity-preserving subalgebra, and at \(L=4\)
exactly one short of \(\dim\mathrm{dTL}_4=M_8=323\).  This shard answers the
two questions those numbers pose: \emph{why} \(322\), and \emph{which}
element spans the kernel.

\subsection{Parity-Refined Multiplicities}

\textbf{Derived.}  Write \(m_c^{\pm}(L)\) for the number of CA-66 basis
paths of \(O^{\otimes L}\) with total charge \(c\in\{\one,\tau\}\) and even
(\(+\)) or odd (\(-\)) occupied number \(N\).  Appending a site keeps \(N\)
(vacancy summand) or increments it and fuses the charge (\(\tau\) summand),
so
\[
  m_\one^{\pm}(L+1)=m_\one^{\pm}(L)+m_\tau^{\mp}(L),\qquad
  m_\tau^{\pm}(L+1)=m_\tau^{\pm}(L)+m_\one^{\mp}(L)+m_\tau^{\mp}(L),
\]
with seed \((m_\one^+,m_\one^-,m_\tau^+,m_\tau^-)(1)=(1,0,0,1)\).  The
alternating sums \(M_c^{-}(L)=m_c^+(L)-m_c^-(L)\) obey the Fibonacci
recursion generated by \(I-T\), \(T=\bigl(\begin{smallmatrix}0&1\\1&1
\end{smallmatrix}\bigr)\), giving \(M_\one^-(L)=F_{L+1}\),
\(M_\tau^-(L)=-F_L\) and the closed forms
\[
  m_\one^{\pm}(L)=\tfrac12\bigl(F_{2L-1}\pm F_{L+1}\bigr),
  \qquad
  m_\tau^{\pm}(L)=\tfrac12\bigl(F_{2L}\mp F_{L}\bigr),
\]
consistent with \(m_\one=F_{2L-1}\), \(m_\tau=F_{2L}\) (CA-66, T66.2).
Charge \(\one\) is even-heavy (the all-vacuum state and the
\(\tau\otimes\tau\to\one\) channel), charge \(\tau\) odd-heavy (single-\(\tau\)
states); the asymmetry \(\sim\varphi^L\) is subdominant to
\(m_c\sim\varphi^{2L}\).
\[
\begin{array}{c|cc|cc}
L & m_\one^+ & m_\one^- & m_\tau^+ & m_\tau^- \\\hline
1 & 1 & 0 & 0 & 1\\
2 & 2 & 0 & 1 & 2\\
3 & 4 & 1 & 3 & 5\\
4 & 9 & 4 & 9 & 12\\
5 & 21 & 13 & 25 & 30\\
6 & 51 & 38 & 68 & 76
\end{array}
\]
The \(L=4\) row is the observed split \(13=9+4\), \(21=9+12\).  The
parity-\emph{preserving} subalgebra
\(\bigoplus_c\bigl(M_{m_c^+}\oplus M_{m_c^-}\bigr)\) has dimension
\[
  P(L)=\sum_c\bigl((m_c^+)^2+(m_c^-)^2\bigr)
  =9,\ 51,\ 322,\ 2135,\ 14445\quad(L=2,\dots,6),
\]
with complement check \(\dim\Alg_L-P(L)=\sum_c2m_c^+m_c^-\)
(\(610-322=288\) at \(L=4\)).  This derives the \(322\): the CA-73 datum
says \(\rho_4\) is \emph{onto} the parity-preserving subalgebra, the
maximum permitted by the parity obstruction of CA-69.

\subsection{The Kernel Theorem}

\textbf{Convention bridge (Rule 3, recorded).}  SRC-TL-JONES works with
\(\mathrm{TL}_n(q)\), loop parameter \(-(q+q^{-1})\), and \(\ell=|q^2|\).
Our diagram algebra has loop value \(\beta=+\varphi=2\cos(\pi/5)\)
(CONVENTIONS (r); \(e_j^2=\varphi e_j\) checked in CA-73), i.e.\
\(q=e^{i\pi/5}\), \(\ell=5\); quantum integers
\([1],\dots,[5]=1,\varphi,\varphi,1,0\).  The loop sign is a gauge choice
(\(\kappa_\tau=+1\) is sourced) and touches none of the dimension, trace,
or radical statements used below; the identification is pinned by the
source's own \(\ell=5\) example, whose module dimensions are exactly our
path counts.  In its notation the first Jones--Wenzl idempotent at the root
is \(E_{\ell-1}=E_4\in\mathrm{TL}_4\); we write \(p_4\).

\textbf{Theorem 74.1 [derived; one numerical input].}  Let
\(\iota:\mathrm{TL}_4(\varphi)\cong\pi_A\,\mathrm{dTL}_4\,\pi_A\) be the
fully-occupied corner isomorphism (\(A=\{1,2,3,4\}\)).  Then
\[
  \ker\rho_4=\Span\{\iota(p_4)\},
\]
and \(\rho_4\) maps \(\mathrm{dTL}_4(\varphi)\) onto the parity-preserving
subalgebra of \(\Alg_4\).  \emph{Proof.}
(i) \emph{\(p_4\) exists and is negligible:} at \(\ell=5\) the trace form
on \(\mathrm{TL}_4\) is degenerate with one-dimensional radical generated
by the unique idempotent \(p_4\) satisfying \(e_ip_4=p_4e_i=0\)
(\(i=1,2,3\)), and \(\operatorname{tr}_4(p_4)=0\) [sourced :63--69,
:135--141, :818].
(ii) \emph{\(\iota(p_4)\neq0\):} \(p_4\) has coefficient \(1\) on the
identity diagram [sourced :792], diagrams are a basis, and \(\iota\) is
injective [sourced corner iso].
(iii) \emph{\(\rho_4(\iota(p_4))=0\):} the corner restriction
\(\rho_4\circ\iota\) is the dense evaluation
\(\mathrm{TL}_4(\varphi)\to\operatorname{End}_{\Cc}(\tau^{\otimes4})\)
(CA-69), whose kernel is exactly the radical \(\Span\{p_4\}\): the quotient
\(Q_4(5)\cong\operatorname{End}(\Delta_q(1)^{\otimes4})\) has dimension
\(2^2+3^2=13=C_4-1\) [sourced :74--80; ledger derived].  Independently:
\(\rho_4\circ\iota\) is a unital \(*\)-representation in the unitary gauge,
so \(\rho_4(\iota(p_4))\) is an orthogonal projection whose quantum trace
is \(\operatorname{tr}_4(p_4)=0\); positivity of quantum dimensions
(CONVENTIONS (r)) forces it to vanish.
(iv) \emph{Rank--nullity:} \(\dim\mathrm{dTL}_4=M_8=323\) [sourced CA-69]
and \(\operatorname{rank}\rho_4=322\) [numerical input, CA-73 datum, re-run
this session], so the kernel is one-dimensional and (ii)+(iii) identify it.
Since \(322=P(4)\) and the image lies in the parity-preserving subalgebra
(CA-69), the map is onto it. \(\qed\)

Every ingredient except \(\operatorname{rank}\rho_4=322\) is analytical and
locally sourced; the numerical rank is the surjectivity half, and the Julia
suite additionally confirms \(\rho_4(\iota(p_4))=0\) to machine precision.

\subsection{The Upgraded Conjecture}

\textbf{Conjecture 74.2 [proposal; refines CA-69].}  For all \(L\ge1\),
\(\rho_L\) maps \(\mathrm{dTL}_L(\varphi)\) onto the parity-preserving
subalgebra of \(\Alg_L\), with kernel the two-sided ideal generated by the
corner Jones--Wenzl embeddings,
\(\ker\rho_L=\langle\iota_A(p_4):A\subseteq[L],\,|A|=4\rangle\).  Exact for
\(L\le3\) (zero kernel) and \(L=4\) (Theorem 74.1; a two-sided ideal inside
a one-dimensional kernel equals it).  Under the conjecture,
\[
  \dim\ker\rho_L=M_{2L}-P(L)=0,\ 0,\ 1,\ \mathbf{53},\ \mathbf{1066}
  \quad(L=2,\dots,6);
\]
because image \(\subseteq\) parity-preserving is \emph{proved}, these are
unconditionally rigorous \emph{lower bounds} on \(\dim\ker\rho_L\) — the
open half at \(L\ge5\) is surjectivity alone.  For comparison, the dense
kernels \(\ker(\mathrm{TL}_n(\varphi)\to
\operatorname{End}(\tau^{\otimes n}))\) have dimensions
\(C_n-F_{2n-1}=1,8,43\) (\(n=4,5,6\)); the dilute kernel is strictly larger
because \(\mathrm{dTL}_L\) glues binomially many dense corners [sourced
:2314--2318, the dilute-radical decomposition].  A within-budget \(L=5\)
check must be combinatorial — either exhibiting dilute words spanning all
\(P(5)=2135\) parity-blocks, or computing
\(\dim\langle\iota_A(p_4)\rangle=53\) through the dilute cellular
structure; the \(2188\)-dimensional rank computation is out of budget
(AGENTS.md platform note).

\subsection{Julia Invariants}

\texttt{src/JonesWenzlKernel.jl} builds the corner unit \(\Pi_A\), the
corner-compressed dense generators \(E_j\), and the Wenzl recursion
\(p_{r+1}=p_r-([r]/[r+1])\,p_rE_rp_r\) on the CA-66 path basis, carrying
the formal word expansion alongside the matrices.  Because the local
extraction of the recursion truncates its diagram factor (source note
above), the testset ``jones-wenzl kernel decision (CA-74)'' pins the
\emph{defining} properties instead of the formula: idempotency,
self-adjointness, two-sided annihilation \(E_jp_k=p_kE_j=0\) for \(j<k\),
and unit-word coefficient exactly \(1\).  It further pins:
\(\operatorname{tr}\Pi_A=(2,3)\) on the two charge blocks; \(e_j^2=\varphi
e_j\) and the braid relation in the corner; \([1..5]=(1,\varphi,\varphi,1,0)\)
with \([5]=0\) at \(10^{-14}\); survival
\(\|\rho(\iota(p_2))\|=\|\rho(\iota(p_3))\|=1\); \textbf{the decision}
\(\|\rho(\iota(p_4))\|_{\mathrm{op}}=2.3\times10^{-15}\le10^{-12}\); the
CA-73 cross-checks \(\texttt{dilute\_image\_dimension}(4)=322=
\texttt{parity\_even\_dimension}(4)\) (reused, not forked); and the full
parity table above for \(L=2..5\) by pure path counting.  113 assertions,
standalone green, re-run independently by the orchestrator; three
mutations (Wenzl coefficient, recursion \(\beta\), corner unit) each RED
(14/25/8 failures), reverted.

\subsection{Gram-Rank Deferral and Boundary}

The Gram-rank kernel dictionary (descendant Gram matrices, JW \(\leftrightarrow\)
null-vector reading) is W2.3 and stays behind the endpoint-closing ledger
W1.5/T1: Gram data need a canonical fusion-tree ONB in each fixed-\(S\)
sector, which no shard yet fixes; CA-74's statement is rank-only and
basis-independent, which is why it could proceed (PLAN.md W1.4).  Not
claimed here: any continuum or Koo--Saleur statement; a closed form for the
dilute ideal \(\langle\iota_A(p_4)\rangle\) at general \(L\); surjectivity
beyond \(L=4\).  Recorded gaps (acquisition queue): a source with the
explicit algebraic Wenzl recursion (e.g.\ Kauffman--Lins or Morrison
arXiv:1503.00384) and the categorical negligible-morphism formalism (e.g.\
EGNO Ch.\ 8 / Goodman--Wenzl); one line of the ILZ extraction
(\texttt{:666}, ``\(\dim Q_8=233\)'') contradicts the paper's own
\(\ell=5\) formulas (\(610\)) and is treated as an extraction artifact —
only \(\dim Q_4(5)=13\) is used, triple-checked.
